\section{Kết luận và hướng phát triển}

\subsection{Chức năng và kết quả đạt được}

Đồ án đã hiện thực trò chơi Bloxorz 3D có thể điều khiển bằng tay, gồm luật cơ
bản, ô Fragile, cầu, ba loại công tắc và cơ chế tách, đổi khối điều khiển, ghép
khối. Bốn bộ giải BFS, DFS, UCS, A* dùng chung \texttt{GameEngine}. Giao diện
hỗ trợ giải trong luồng nền, kiểm chứng và phát lại lời giải. Bộ kiểm thử đạt
203/203 trường hợp; cả 180 lượt đo hiệu năng đều tìm
thấy lời giải và tập lệnh độc lập xác nhận 60/60 lời giải hợp lệ.

Trên 15 màn chơi, BFS, UCS và A* có cùng số bước di chuyển trung bình. Khi
đánh giá lại bằng hàm chi phí chuẩn, đường BFS có chi phí bằng UCS trên cả 15
màn chơi. DFS có thể nhanh nhưng lời giải dài và nhạy với thứ tự sinh trạng
thái kế tiếp, nổi bật là 724 bước ở màn chơi 15. A* giảm trung bình 43,4\% số
nút mở rộng so với BFS và 47,8\% so với UCS, nhưng không luôn nhanh hơn: ở màn
chơi 15, A* mất 631,90 ms so với
534,95 ms của BFS.

\subsection{Giới hạn của kết luận}

Heuristic A* dùng cận dưới dựa vào lần dịch chuyển tức thời cuối. Lập luận về
tính chấp nhận được với hàm chi phí chuẩn nằm ở Mục~\ref{sec:algorithms}.
Bảo đảm tối ưu của A* không tự áp dụng cho hàm chi phí hoặc heuristic tùy chỉnh.
Các bảo đảm tìm kiếm cũng đòi hỏi thuật toán không bị hủy, hết thời gian hoặc chạm giới
hạn nút. Cận hiện hành có thể hướng dẫn tìm kiếm kém hiệu quả ở màn chơi có
nhiều lần dịch chuyển tức thời.

Các nhận xét thực nghiệm chỉ dựa trên 15 màn chơi và ba lần lặp. Bộ nhớ đỉnh
là lượng cấp phát Python do \texttt{tracemalloc} theo dõi, không phải tổng RAM
của tiến trình; thời gian trên màn chơi nhỏ có thể chịu nhiễu hệ điều hành.

\subsection{Hướng phát triển}

Có thể thiết kế cận dưới mạnh hơn có xét cấu trúc kết nối của cầu và công tắc,
rồi chứng minh và đo lại trước khi thay heuristic mặc định. Các thử nghiệm tiếp
theo nên thêm màn chơi lớn, màn chơi sinh tự động, nhiều lần lặp với thứ tự
ngẫu nhiên, đo RSS và hiển thị sai số trên biểu đồ. Về phần mềm, có thể bổ sung
trình biên tập màn chơi, lưu tiến độ và hoạt ảnh phong phú hơn ở tầng giao diện.
